\section{5교시 (75분) — 교훈과 편익: 프로젝트가 끝난 뒤에 남는 것}

% =====================================================================
% 5교시 — 교훈과 편익 (교재 제5장, 워크북 §4~§5)
% =====================================================================
\begin{frame}{5교시에서 배우는 것 — 프로젝트가 사라진 뒤의 독자를 위해 쓴다}
\begin{itemize}
\item 교훈이 조직에 남지 않는 이유 — 수집이 아니라 \textbf{흡수}의 실패 (Williams · Syllk)
\item 교훈 등록부 6칸 — 적용 대상·소유자 없는 교훈은 감상 / 온담 15건 \ra{} 개정 요청 3건
\item 편익의 책무 분리(ITO) — output은 G4에서 종결, outcome은 2029-05까지 실명 서명
\item 분모 정치 91.3 vs 84.6 — 승인된 목표는 승인된 분모로 판정한다
\item 손익분기 66\%·두 시계·편익 리스크 5건 — 그리고 P1은 어느 성공인가
\end{itemize}
\keymsg{교훈과 편익의 공통 실패는 하나 — 프로젝트가 사라진 뒤의 독자를 위해 쓰지 않은 것이다}
\note{제5장 도입. 4교시까지가 "종료를 증명하는 문서"였다면 5교시는 "종료 뒤에도 살아 있어야 하는 문서" 둘 — ⑲ 교훈 등록부와 ⑳ 편익 실현 원장 — 을 다룬다. 둘의 독자는 P1 PM이 아니라 06-01 이후의 프로그램 PMO장·재경팀장·오정근 본부장이며, 그래서 모든 칸에 PM이 아닌 실명이 들어간다. 데이터 일자 2027-05-28(G4), 원장 v1.0 서명 05-21, 개정 요청 발송 06-04. 워크북 §4·§5 완성 예시본이 정본이고 교재와 다르면 워크북을 따른다. 예상 소요 4분.}
\end{frame}

\begin{frame}{교훈은 왜 조직에 남지 않는가 — Williams(2008)·Syllk의 여섯 층}
\fig[0.6]{../그림/PM_Day4/fig_5_1_syllk.pdf}
\figcap{그림 5-1. Syllk 6층(사람 3 / 시스템 3)과 교훈 6칸. ④ 적용 대상·⑤ 소유자가 사람 층에서 시스템 층으로 가는 문이고, 개정 요청서가 그 문을 통과하는 유일한 형식 — 6개월 소멸선은 [추가 설정]}
\keymsg{교훈이 남는 조건은 교훈의 수가 아니라 — 그 교훈으로 개정된 양식이 다음 프로젝트의 착수 패키지에 있는가다}
\note{교재 5.1절. Williams(2008, PMJ 39(4))의 실증 — 교훈 학습의 실패는 모으지 못해서가 아니라 모은 교훈이 절차(양식·규정·계약 문안)로 흡수되지 않아서다. Duffield \& Whitty(2015, IJPM)의 Syllk 모형은 흡수 조건을 여섯 층으로 구조화한다 — 사람 쪽 learning·culture·social, 시스템 쪽 technology·process·infrastructure. 회고 워크숍과 PM 커뮤니티는 사람 층만 통과시키고, 시스템 층에 닿지 않으면 6개월 뒤 사라진다. 이것은 Day1 1.3절 임시조직 이론과 같은 말이다 — 임시조직이 해산하면 배운 것도 해산하므로 해산 전에 영구조직의 제도로 옮겨야 한다. 온담이 15건마다 적용 대상을 세 표준 개정 요청으로 묶는 이유가 여기 있고, 개정 요청서는 등록부의 부속이 아니라 결과물이다. 틀리는 방식 셋 — 저장소에 넣고 끝, "검토 바람", 회고를 의식으로. 예상 소요 6분.}
\end{frame}

\begin{frame}{교훈 등록부의 여섯 칸 — "적용 대상"이 없는 교훈은 감상이다}
\begin{tbl}[\scriptsize]{L{0.14\textwidth}L{0.42\textwidth}L{0.36\textwidth}}
\textbf{칸} & \textbf{무엇을 적는가} & \textbf{규율} \\ \midrule
① 상황 & 무엇이 있었나 — 사실·일자·숫자 & 형용사 없이 \\
② 원인 & 왜 그랬나 — 양식·규칙·구조 & \textbf{사람 이름 금지} \\
③ 교훈 & 무엇을 바꾸면 되나 — 한 문장 & "중요하다"는 감상 \\
④ 적용 대상 & 어느 양식·규칙·문안의 어느 칸 (NP / STD / CTR) & "전사"는 없음과 같다 \\
⑤ 소유자 & 누가 바꾸나 — 실명 & PM은 06-01에 없어진다 \\
⑥ 증빙 & 어느 문서의 어느 행 & 워크북 ⑫\til{}⑯ 항목 번호 \\
\end{tbl}
\begin{itemize}
\item 잘된 것도 교훈이다 — 우연한 성공을 양식으로 고정하지 않으면 다음 프로젝트가 다시 발명해야 한다
\end{itemize}
\keymsg{"담당자가 놓쳤다"는 사실이지 교훈이 아니다 — "양식에 그 칸이 없었다"로 적어야 칸이 다음 담당자를 잡는다}
\note{교재 5.2절, 워크북 §4 항목 1. 여섯 칸이 다 채워질 때만 교훈이고 하나라도 비면 감상이다. 원인 칸에 사람을 적지 않는 것이 이 설계의 규율 — 다음 프로젝트의 담당자도 놓칠 것이므로, 양식이 바뀌어야 그 뒤의 모든 프로젝트에서 칸이 잡는다. 사람 이름은 소유자 칸에만 나온다. 적용 대상 약호 NP 다음 프로젝트 양식 / STD 사내 표준 / CTR 계약 문안. 온담 Lessons log 41건 중 종료 정리 15건(나머지 26: CIR 12·P3/P4 소관 9·재발 낮음 5), 원천 분포 조기 신호 5 / 착시 1 / 감리·계약 3 / 이슈·리스크 2 / 절차·거버넌스 3 / 잘됨 1(L-15; L-08도 잘됨). 검산 15 = 5 + 1 + 3 + 2 + 3 + 1. 예상 소요 6분.}
\end{frame}

\begin{frame}{온담 P1 교훈 15건 중 다섯 — 원인에 사람이 없고 소유자에 PM이 없다}
\begin{tbl}[\scriptsize]{L{0.07\textwidth}L{0.24\textwidth}L{0.30\textwidth}L{0.17\textwidth}L{0.12\textwidth}}
\textbf{ID} & \textbf{상황} & \textbf{원인(구조) \ra{} 교훈} & \textbf{적용 대상} & \textbf{소유자} \\ \midrule
L-02 & 승인 대기 12영업일, 대기 비용 0.42 & 헌장·RACI에 대체 결정권자 칸 없음 \ra{} 실명 지정 & STD ① 헌장 §11 & PMO장 \\
L-06 & 감리 AF-1-03·04 귀속 다툼, 0.18 & 감리 계약에 귀속 판정 기준 없음 \ra{} WBS 담당 열·RACI & CTR 감리 부속서 3 & 강현도·정하늘 \\
L-08 잘됨 & 대기 비용 판정 1회로 종결 & 회의록 "요청·응답 일자" 열 \ra{} 동시 기록 서명 & CTR §14 / NP 회의록 & 강현도·P1 PM \\
L-09 & D$-$2 라이선스 CFO 결재 대기, 0.02 & 창 앞 사전 결재 규정 없음 \ra{} 4주·5천만 원 일괄 결재 & STD ② 내부통제 & 재경팀장·정미란 \\
L-15 잘됨 & 두 컷오버 정시·트리거 0 & 1차 지표를 2차 go 조건으로 \ra{} 컷오버 표준 & STD ① 컷오버 표준 & PMO장·박준영 \\
\end{tbl}
\keymsg{L-02의 원인을 "본부장이 바빴다"로 적으면 사실이지만 교훈이 아니다 — 다음 프로젝트의 Senior User도 바쁠 것이다}
\note{교재 5.3절, 나머지 10건은 워크북 §4.5. 다섯 건은 Day3·Day4의 사건이 계약 문안과 사내 표준으로 바뀌는 경로를 보여 준다 — L-02는 Day3 ⑬ IS-02·⑫ 신호 2, L-06은 Day3 4.3절 "등급보다 귀속 칸이 무겁다"가 계약 부속서가 되는 길, L-08은 동시 기록(contemporaneous records) 덕에 4.3절의 지체상금 0이 네 줄로 증명된 근거, L-09는 2.8절 D$-$2 사건, L-15는 R-14 트리거를 게이트 조건으로 묶은 사슬 — 어느 문서에도 "표준"으로 적혀 있지 않았다. L-08의 소유자에 P1 PM이 있는 것은 06-01에 프로그램 PMO장으로 승계된다는 뜻이다. 실무에서 틀리는 방식 다섯 — 원인에 사람, 적용 대상 "전사", 소유자 PM, 잘된 것 없음, 개정 요청이 부속. 예상 소요 7분.}
\end{frame}

\begin{frame}{사내 표준 개정 요청 3건(06-04 발송) — 15건을 15곳에 보내면 15건이 "검토 중"이 된다}
\begin{tbl}[\scriptsize]{L{0.22\textwidth}L{0.17\textwidth}L{0.12\textwidth}L{0.41\textwidth}}
\textbf{표준} & \textbf{소유자} & \textbf{의결 시점} & \textbf{개정 조항(교훈)} \\ \midrule
① 프로젝트 관리 규정·양식집 & 프로그램 PMO장 & 2027-07 운영위 & 성과보고 SV·CV 분해 열(L-01·03) · 헌장 대체 결정권자(L-02·10) · RoC 부속(L-11) · 컷오버 표준(L-15) · BC 편익 표 분모·실명 필수(⑳) \\
② 재무본부 내부통제 지침 & 재경팀장·정미란 & 2027-08 & 5,000만 원 조항 단서 신설 · 개인정보 사건 판단 시한 절차(L-09·14) \\
③ 구매 계약 표준 문안 & 강현도·최영주 & 2027-09 P2 후속 계약부터 & SI 계약 §14 여유 소유권·동시 지연·대기 기록 · 감리 부속서 3 귀속 기준 · 조달 상태 양식(L-06\til{}09) \\
\end{tbl}
\begin{itemize}
\item 조항 번호와 새 문장이 있는 요청만 의결된다 — "검토 바람"은 영원히 검토된다
\item 적용 시점이 P2 G2 전 — 다음 프로젝트의 착수 패키지에 들어가야 흡수다
\end{itemize}
\keymsg{개정 요청서는 교훈 등록부의 부속이 아니라 결과물이다 — 아카이브에 함께 들어가면 함께 잊힌다}
\note{교재 5.3절 후반, 워크북 §4 항목 2. 세 표준은 각각 소유 부서와 개정 절차가 하나이므로 15건을 세 묶음으로 보내면 세 번의 의결로 끝나고, 15곳에 보내면 15건이 "검토 중"이 된다. 각 요청서에는 현행 조항 번호·새 문장·근거 교훈 ID·적용 시점이 있고, ①은 P2 G2 전, ③은 P2 후속 계약부터 적용된다. 지속적 성공(Cooke-Davies)의 판정 기준이 바로 이 세 개정이 07·08·09월에 실제 의결되어 P2·P5·2028 후속 과제 착수 패키지에 들어갔는가이다 — 마지막 프레임에서 회수한다. 예상 소요 5분.}
\end{frame}

\begin{frame}{편익 실현의 책무 분리 — Zwikael \& Smyrk의 ITO, 측정 책임자 실명이 서명이다}
\fig[0.58]{../그림/PM_Day4/fig_5_2_ito_ledger.pdf}
\figcap{그림 5-2. Input \ra{} Transform \ra{} Output $|$ Utilisation \ra{} Outcome과 책임 분리선. output 서명(P1 PM·이재현)은 G4 인수로 종결, outcome(오정근 5·한동석 1)·utilisation 서명은 2029-05까지 유효(BR-04) — 원장 6칸의 ⑤ 칸이 서명란}
\keymsg{P1 PM은 폐기율을 내릴 수 없다 — PM이 outcome에 서명하면 없어진 조직이 책임을 지고, 그것은 아무도 지지 않는 것과 같다}
\note{교재 5.4절, 워크북 §5 항목 5. Day1 3장의 ITO 모형 회수 — output(플랫폼)의 책임은 프로젝트에, outcome(폐기율 2.2\%)의 책임은 프로젝트 소유자(현업 책임자, PRINCE2 Senior User)에, 그 사이 utilisation(매장이 발주 가이드를 지키고 구매팀이 폐기를 집계하는 것)이 있다. 이 분리가 없으면 G5에서 편익 미달의 원인은 "P1이 잘못 만들었다"와 "현업이 안 썼다" 사이에서 영원히 다투어진다. ITO 서명표 05-21: output P1 PM·이재현 / outcome 오정근(B-04·05·07·03·08)·한동석(B-02) / utilisation 매장운영팀장·강현도·마케팅·인사본부·실사 담당·나영선·P3 PM·박준영·재경팀장 / 판정 재무본부(정미란·재경팀장). 서명자 변경 시 후임이 재서명한다. Day1 BC §4 "측정 책임자를 부서명이 아니라 실명으로"가 이 날을 위한 것이었다. 반대 실패 — outcome 소유자가 서명하지 않으면 "IT 시스템의 성과"가 되어 현업은 쓰지 않아도 면책된다. 예상 소요 6분.}
\end{frame}

\begin{frame}{편익 실현 원장 6칸 — 직접 편익 118 / 온담 귀속 99.6억}
\begin{tbl}[\scriptsize]{L{0.07\textwidth}L{0.30\textwidth}L{0.24\textwidth}L{0.15\textwidth}L{0.13\textwidth}}
\textbf{B-ID} & \textbf{편익(측정식·리포트)} & \textbf{기준값 \ra{} 목표} & \textbf{측정 책임자} & \textbf{승인 / 귀속} \\ \midrule
B-04 & 폐기율 = 월 폐기 금액 $\div$ 취급원가 (RPT-INV-05) & 3.9\%(2025) \ra{} 2.2\%(M+24) & \textbf{오정근} / 강현도 & 43 / 43 \\
B-05 & 배달수수료 회피 = 전환 매출 $\times$ 12.8\% (RPT-ORD-02) & 앱 2.1\% \ra{} 온라인 18\% & \textbf{오정근} / 재경팀장 & 37 / \textbf{18.6} \\
B-07 & 직영 인건비율(분기) & 21.4\% \ra{} 20.1\% & \textbf{오정근} / 인사본부 & 38 / 38 \\
B-02 & 재고 정확도, 두 분모 병기 (StRS-034) & A 91.3 \ra{} 98.5 / B 84.6 미설정 & \textbf{한동석} / 실사 담당 & (B-04 기여) \\
B-08 & 가맹 발주 리드타임 (RPT-FRN-03) & 38h \ra{} 12h · \textbf{첫 측정 13.6h} & \textbf{오정근} / 박준영 & (B-03 기여) \\
\end{tbl}
\begin{itemize}
\item 6칸 = 측정식(코드 수준) · 기준값(값·시점·분모·출처) · 목표값 · 측정 시점 · 측정 책임자 실명 · 미달 시 조치
\item 미달 시 조치는 "원인 분석"이 아니라 output / outcome / utilisation 세 갈래 중 하나로 분류 \ra{} 그 갈래의 소유자가 움직인다
\end{itemize}
\keymsg{원장 v1.0 2027-05-21 서명, 정본 관리자는 재무본부 — P1 종료 후에도 살아 있어야 하는 유일한 P1 문서다}
\note{교재 5.4절 표, 워크북 §5 항목 1. B-03 결품률(4.7 \ra{} 1.5\%)은 기회손실 산정식이 2027-11 확정이라 금액 미산정으로 표에서 뺐다. 검산 118 = 43 + 37 + 38, 온담 귀속 99.6 = 43 + 18.6 + 38 — B-05의 37은 이사회 승인본 숫자이지만 가맹 매장의 배달 매출은 가맹점주의 것이므로(C-13은 연결매출이 아니다) 온담 손익 귀속은 직영 비중 2,980/5,920 = 50.3\%인 18.6억, 나머지 18.4억은 가맹점주 편익이다. B-08 진행률 94\% = (38 $-$ 13.6)/(38 $-$ 12). 미달 시 조치의 분기를 B-04로 예시 — output(리포트 오류) \ra{} 하자보수 / outcome(발주 정확도·수요 예측 미사용) \ra{} 오정근의 발주 가이드 채택 조치 / utilisation(매장 미준수) \ra{} P3 교육·지역매니저 점검. 미달 판정 시 오정근이 G5 전 개선 계획을 낸다. 예상 소요 6분.}
\end{frame}

\begin{frame}{분모 정치 — 91.3 vs 84.6, 승인된 목표는 승인된 분모로 판정한다}
\fig[0.58]{../그림/PM_Day4/fig_5_3_denominator.pdf}
\figcap{그림 5-3. A 순환실사 SKU(69\%) 91.3 \ra{} 98.5(7.2\%p) vs B 전체 SKU 84.6(13.9\%p, 1.93배). 비순환 31\%의 검사기록은 종이·엑셀(2025-11 이사회 반려 34억) — 판정 A 서명 05-21, B 목표는 이관 6번 가결/부결(BR-05) 뒤 운영 +12에서 설정}
\keymsg{Day3의 정정(R-12)은 실은 편익 판정권의 문제였다 — 누가 어느 분모로 98.5를 판정하는가를 서명으로 확정하지 않으면 G5는 협상이 된다}
\note{교재 5.6절, 워크북 §5 항목 2. 분모 A = 순환실사 SKU 기준 91.3\%(물류본부 2025 보고, O-08), 분모 B = 전체 SKU 기준 84.6\%(2026-12-18 StRS-034 두 분모 리포트 첫 실행 — Day2에서 Must로 둔 이유). 차이의 원인은 비순환 SKU(원료·반제품·포장재)의 입고 검수·검사기록이 종이·엑셀이라 실사 정확도 자체가 낮고, P1은 그 디지털화를 하지 않았다(헌장 §5 제외 범위). 논리 셋 — ① 이사회 승인본의 분모는 "순환실사 SKU"로 명기 ② B로 바꾸면 개선 폭 13.9\%p, 약속하지 않은 것을 약속한 셈 ③ B의 정확도는 P1 밖의 전제 위에 있다. 합의 한동석·나영선·재경팀장 05-21, 이사회 06월 보고에 "두 분모 병기·판정은 A" 명시. 유리한 분모를 택한 것이 아니라 약속한 분모로 판정하고, 약속하지 않은 분모의 전제(이관 6번, 나영선 2027-09 부의)를 정직하게 적은 것이다. 틀리는 방식 넷 — 분모를 안 적음, 유리한 분모 사후 선택, 불리한 분모로 "정직" 과시(전제 미충족이 P1 결함으로 기록됨), 전제 미이관. 예상 소요 7분.}
\end{frame}

\begin{frame}{편익이 실패하는 이유 — Ward \& Daniel의 셋과 온담의 세 가지 정치}
\begin{columns}[T]
\begin{column}{0.48\textwidth}
\subhead{실패 원인 셋 (Ward \& Daniel 2012)}
\begin{itemize}
\item 소유자 부재 — "회사"의 편익은 아무의 것도 아니다
\item 측정 시점 부재 — M+24가 달력의 어느 날인가
\item 기준값의 정치 — 사후에 정해진 기준값은 협상이 된다
\end{itemize}
\end{column}
\begin{column}{0.50\textwidth}
\subhead{온담 원장의 세 정치}
\begin{itemize}
\item 분모 정치 — A 91.3 / B 84.6, 판정은 A
\item 귀속 정치 — B-05 37 \ra{} 온담 18.6 + 가맹점주 18.4
\item 측정 조직 소멸 — B-08 로그: P1이 만들고 P5가 읽고 영업본부가 해석
\end{itemize}
\end{column}
\end{columns}
\keymsg{기준값은 프로젝트 전에, 측정 책임자는 프로젝트 중에, 판정은 프로젝트 후에 — 셋이 같은 문서에 있어야 한다}
\note{교재 5.5절, 워크북 §5.1. 편익 의존성 네트워크 — IT 산출물과 편익 사이에 반드시 업무 변화가 있고 그 변화의 소유자가 편익의 소유자다(ITO의 utilisation을 다른 방향에서 본 것). 귀속 정치를 적어 두지 않으면 G5에서 재무본부가 손익을 대조하는 순간 "37억이 어디 있는가"가 되고 그때 처음 설명하면 P1 편익 전체가 의심받는다. 적어 두면 18.4억은 협의회(서정필)에 "가맹점주 편익"으로 공유할 수 있는 숫자가 되어 부속서 데이터 용도 논쟁(R-02·BR-03)의 지렛대가 된다. 측정 조직 소멸은 세 조직 중 하나가 바뀌면 측정이 끊기는 리스크 — BR-04로 로그 보존 3년·서명자 변경 시 재서명·정본 재무본부. 틀리는 방식 — 승인 금액 전액 귀속, "M+24"만 적고 달력 날짜·시계 없음, 로그 보존 정책 없음. 예상 소요 5분.}
\end{frame}

\begin{frame}{손익분기 66\% · 62.1\% · 78\% — 그리고 M+12 / M+24 리뷰의 두 시계}
\fig[0.56]{../그림/PM_Day4/fig_5_4_breakeven_clocks.pdf}
\figcap{그림 5-4. NPV 대 실현률(50\% $-$58.3 / 66\% 0 / 70\% +14.8 / 100\% +124.4; AC 151.95 반영 시 손익분기 62.1\%·+138.3), 판정선 95\%·66\%, 온담 귀속 기준 78\% + 프로그램 시계(M+24 2028-01·G5 2028-03-31)와 운영 시계(+12 2028-05·+24 2029-05)}
\keymsg{BC의 시계는 고치지 않는다 — 승인 문서를 사후에 고치는 것은 기준값을 고치는 것과 같다. 대신 운영 시계를 추가한다}
\note{교재 5.7절, 워크북 §5 항목 3·4. 손익분기 66.0\% = 118 $\times$ 0.66 = 77.8억/년(할인율 8\%, 운영비 27억 전액 P1 귀속, 보수적). 최종 AC 151.95 반영 시 62.1\% — 컨틴전시 반납 4.05·관리예비비 12.0 미사용의 효과가 3.9\%p 하락이다. 온담 귀속 99.6 기준으로는 요구 실현률이 78\%(77.8/99.6)로 올라간다 — 이것도 원장에 적혀야 G5가 정직하다. 판정 규칙(허용범위 $-$5\%p): 95\% 미만 경고 / 66\% 미만은 헌장 §14 중단 조건 수준 — 프로젝트는 종료되었으므로 프로그램 이사회의 후속 투자 판단 근거. 두 시계 — BC는 프로그램 착수 2026-01 기준이라 B-08의 M+12가 컷오버 전(2027-01)이 되는 모순이 있지만 고치지 않고, 운영 이관 2027-05 기준 시계를 추가한다(운영 +12 원장 v2: B-02 분모 B 목표 설정·측정 조직 확인, +24 원장 v3: 유지 판정·원장 종결 또는 상시 KPI 전환). 첫 측정 2027-04 B-08 / 06 B-04·02·03 / 09 B-05·07 / 11 B-03 산정식. 예상 소요 6분.}
\end{frame}

\begin{frame}{편익 리스크 5건 — 원장의 부속, 자금원은 운영 예산·프로그램 예비비}
\begin{tbl}[\scriptsize]{L{0.08\textwidth}L{0.26\textwidth}L{0.06\textwidth}L{0.48\textwidth}}
\textbf{ID} & \textbf{리스크} & \textbf{P} & \textbf{대응} \\ \midrule
BR-01 & 분모 분쟁 재점화(R-12 이관) & 30\% & 05-21 합의문·이사회 보고에 "판정은 A" 명시 \\
BR-02 & 외생변수 — 원재료가·이상 기후로 폐기율 왜곡 & 40\% & 수량 기준 병기, 원재료가 지수 통제 \\
BR-03 & B-05 귀속 분쟁 — 협의회가 18.4억을 본사 편익으로 오해 & 25\% & 귀속 분리 명기·가맹 귀속 지표 공유 \\
BR-04 & 측정 조직 소멸 — 로그·읽는 조직·해석 조직 셋 중 하나 변경 & 35\% & 로그 보존 3년·서명자 변경 시 재서명·정본 재무본부 \\
BR-05 & 전제 미충족 — 검사기록 디지털화 부결 시 B 목표 불가 & \textbf{50\%} & 이관 6번 부의안에 B-02·B-04 연결 명시 \\
\end{tbl}
\begin{itemize}
\item Day1 프리모템 14번 "측정 책임자가 경기 탓이라 했다" \ra{} BR-02·BR-04로 대응된 것
\end{itemize}
\keymsg{편익 리스크는 프로젝트 리스크와 다르다 — 프로젝트는 끝났고, 리스크의 소유자는 원장에 서명한 사람들이다}
\note{교재 5.7절 후반, 워크북 §5 항목 6. 편익 리스크 5건은 원장 부속이며 자금원이 프로젝트 예비비가 아니라 운영 예산·프로그램 예비비다 — P1 예비비는 G4에서 반납되었다(4.05). BR-05의 확률이 가장 높은 이유는 2025-11 이사회가 검사기록 디지털화 34억 SI 견적을 이미 한 번 반려했기 때문이고, 부결 시 B 목표는 "A와 동일 방법으로 측정 가능한 범위"로 한정하며 B-04의 원료 폐기 부분도 미실현으로 남는다. 귀인(BR-01·03)·외생(BR-02)·측정 조직 소멸(BR-04)·전제(BR-05)의 네 유형은 여러분 회사의 편익 원장에도 그대로 적용된다. 예상 소요 4분.}
\end{frame}

\begin{frame}{2027-05-28의 P1은 어느 성공인가 — Cooke-Davies의 셋과 Flyvbjerg의 곱셈}
\begin{tbl}[\scriptsize]{L{0.2\textwidth}L{0.15\textwidth}L{0.55\textwidth}}
\textbf{성공(Cooke-Davies 2002)} & \textbf{판정} & \textbf{근거} \\ \midrule
프로젝트 관리 성공 & \textbf{허용범위 이내} & 재기준선 0 · 컷오버 2회 정시 · VAC $-$5.13(한도 내, PMB 대비 3.5\% 초과) · FP 12,585 100\% · 밀도 0.094 \\
프로젝트 성공(편익) & \textbf{판정 불가} & 첫 측정 하나(B-08 13.6h, 94\%)뿐 — 판정은 G5 2028-03-31 \\
지속적 성공 & 더 뒤 & 세 표준 개정(07·08·09월)이 P2·P5 착수 패키지에 들어갔는가 \\
\end{tbl}
\begin{itemize}
\item "온버짓"이라는 말을 세 층(승인·한도·PMB) 중 어느 것으로 쓰느냐에 따라 답이 다르다 — 이 교재는 그 말을 쓰지 않는다
\item Flyvbjerg의 곱셈 — 원가·일정 두 인수는 확정, 편익 인수는 아직 곱해지지 않았다
\end{itemize}
\keymsg{종료 보고서의 편익 축 "미판정 — ⑳ 원장으로 이관"이 정직한 문장이다 — 원장·서명·두 시계·리스크는 장치이지 결과가 아니다}
\note{교재 5.8절. Day1 1.2절의 두 개념 회수 — 세 성공은 서로 독립이고, Flyvbjerg(2014)의 곱셈은 그 독립성의 통계적 표현이다. 2019년 스마트오더가 Day1 그림 1-6에서 "온타임·온버짓 + 편익 없음" 칸에 놓였던 것을 기억시켜라 — P1이 그 칸에 놓이지 않으리라는 보장은 2027-05-28 시점에 없다. 종료 보고서에 "편익 달성 예정"을 적는 것이 가장 흔한 실패이며, 그러면 G5는 판정이 아니라 협상이 되고 2020년 지적이 2028년에 되풀이된다. 5.9절 실무 실패의 공통점 — 교훈에서도 편익에서도 프로젝트가 사라진 뒤의 독자를 위해 쓰지 않은 것. 생각해 볼 질문 3·5번(여러분이 종료한 프로젝트의 편익에 여러분이 서명했는가 / P1을 "성공"이라 부를 수 있는가)을 짝과 2분 토론. 예상 소요 6분.}
\end{frame}

\begin{frame}{5교시 핵심 정리 — 교훈은 양식으로, 편익은 실명으로 남는다}
\begin{itemize}
\item 교훈 실패는 흡수의 실패 — 시스템 층에 닿는 문은 적용 대상·소유자 칸, 형식은 개정 요청서
\item 6칸이 다 차야 교훈, 원인에 사람이 있으면 교훈이 아니다 — 15건 \ra{} 3표준(07·08·09월)
\item 편익 책무는 output(G4 종결) / outcome(오정근·한동석) / utilisation — 실명 칸이 서명
\item 승인된 목표는 승인된 분모로 — A 91.3 \ra{} 98.5 판정, B 84.6은 전제 결정 뒤 운영 +12
\item 손익분기 66.0 / 62.1 / 78\% · 두 시계 · BR 5건 — P1은 관리 성공은 이루고 편익 성공은 미판정
\end{itemize}
\keymsg{프로젝트는 06-01에 없어진다 — 그 뒤에 남는 것은 개정된 양식 세 개와 서명된 원장 하나다}
\note{제5장 핵심 정리 다섯 줄. 6교시 실습 ②(워크북 §4~§5)로 연결 — ⑲ 교훈 등록부 15건 중 조기 신호 5건을 여섯 칸으로 채우고 개정 요청서 ① 조항 하나를 문장으로 쓰는 것, ⑳ 원장에서 B-02 분모 합의문과 B-05 귀속 분리를 적는 것이 과제다. 검산을 미리 알려라 — 15 = 5 + 1 + 3 + 2 + 3 + 1, 118 = 43 + 37 + 38, 99.6 = 43 + 18.6 + 38, 78\% = 77.8/99.6, B-08 진행률 94\%. §4.7 기본 과제(G2 중단 시 premature closure의 교훈 10건)는 1.2절과 이 장의 결합, §5.7 기본 과제(G5 첫 12개월 실적으로 미달 시 조치 실행)는 5.4\til{}5.7절의 결합 문제. 예상 소요 3분.}
\end{frame}
